% ===================================================================
%  Семинар 6. Содержательная часть — единый источник для обоих файлов.
%
%  Состав и порядок задач — по «Семинар 6.docx» («Гетероскедастичность»):
%      Задание 1 — задача 6.5 (двухшаговая процедура при σ_t^2 = μ + δx_t);
%      Задание 2 — задача 6.6 («реальная задача + Дарбин–Уотсон»);
%      Задание 3 — задача 6.7 («также задача + разные тесты»).
%  Источник: Сборник задач к начальному курсу эконометрики (РАНХиГС),
%  гл. 6 «Гетероскедастичность и корреляция по времени», с. 179–183;
%  условия — по учебнику Магнуса–Катышева–Пересецкого.
%
%  Данные задачи 6.7 — «Семинар 6 (приложение).xlsx» (там же вывод Пакета
%  анализа для регрессии C на Y). Тесты в 6.7 — те, что в сборнике
%  (график, Голдфелд–Квандт, Бройш–Паган), плюс тест Уайта, тест Глейзера
%  (регрессии модулей остатков на функции дохода)
%  и ранговая корреляция Спирмена («Ранг Y» в приложении).
%
%  ОБОЗНАЧЕНИЯ. Как во всех семинарах: RSS — объяснённая сумма квадратов,
%  ESS — остаточная. В сборнике в тесте Бройша–Пагана «RSS/2» — тоже
%  объяснённая сумма вспомогательной регрессии, то есть совпадает с нашей.
%
%  Раздел «Основные формулы» и всё внутри \begin{solution}...\end{solution}
%  попадает только в seminar_06_solutions.pdf. Блоки \begin{excel}...\end{excel}
%  печатаются в обоих файлах.
%
%  Теория сверена с «Теория эконометрики_1_презентация.pdf»:
%  МНК — слайды 10–14, 18–22, 26–27; гетероскедастичность — 42–52;
%  перечень тестов (включая Глейзера и Спирмена) — 48; DW и AR(1) — 60–68.
%  В 6.6 используется Кохрейн–Оркатт (слайд 68); стандартные ошибки
%  берутся из конечной преобразованной регрессии при фиксированном rho_hat.
% ===================================================================

\seminartitle{Семинар 6}{Гетероскедастичность и автокорреляция ошибок}

% ===================================================================
\ifsolutions
% Editable vector illustrations for seminar 6 (TikZ).
% Coordinates of the residual plots (problems 6.6 and 6.7), computed from the task data.
\newcommand{\olsTimeCoordinates}{(1.000000,-0.070722) (2.000000,-0.411736) (3.000000,-0.419468) (4.000000,-0.106693) (5.000000,-0.637064) (6.000000,-0.136421) (7.000000,0.065031) (8.000000,-0.022058) (9.000000,0.143031) (10.000000,0.031247) (11.000000,0.121186) (12.000000,0.179866) (13.000000,-0.004311) (14.000000,0.066977) (15.000000,-0.134230) (16.000000,0.051990) (17.000000,0.056645) (18.000000,-0.082867) (19.000000,-0.061657) (20.000000,0.277906) (21.000000,0.099646) (22.000000,0.571734) (23.000000,0.139581) (24.000000,0.282384)}
\newcommand{\arTimeCoordinates}{(2.000000,-0.399915) (3.000000,-0.264760) (4.000000,0.066042) (5.000000,-0.612357) (6.000000,0.196608) (7.000000,0.140855) (8.000000,-0.066946) (9.000000,0.115915) (10.000000,-0.025101) (11.000000,0.077717) (12.000000,0.131712) (13.000000,-0.080380) (14.000000,0.081664) (15.000000,-0.203237) (16.000000,0.141121) (17.000000,0.012521) (18.000000,-0.093095) (19.000000,-0.027291) (20.000000,0.262160) (21.000000,0.007329) (22.000000,0.471971) (23.000000,-0.106096) (24.000000,0.173564)}
\newcommand{\olsIncomeCoordinates}{(0.000000,-0.680448),(0.000000,-0.188018),(0.000000,0.550625),(1.000000,-0.884880),(1.000000,-0.146237),(1.000000,0.838622),(2.000000,-0.596884),(2.000000,0.141760),(2.000000,1.619048),(3.000000,-0.801317),(3.000000,-0.062673),(3.000000,0.675971),(4.000000,-0.759535),(4.000000,-0.267106),(4.000000,0.225324),(5.000000,-1.210182),(5.000000,0.020891),(5.000000,1.005749),(6.000000,-1.660830),(6.000000,0.062673),(6.000000,1.786175),(7.000000,-1.372833),(7.000000,0.104455),(7.000000,1.089313),(8.000000,-0.838622),(8.000000,0.638666),(8.000000,2.115954),(9.000000,-2.027913),(9.000000,-0.550625),(9.000000,1.172877)}
\newcommand{\wlsIncomeCoordinates}{(0.000000,-0.960880),(0.000000,-0.250066),(0.000000,0.816154),(1.000000,-1.164000),(1.000000,-0.179797),(1.000000,1.132475),(2.000000,-0.728835),(2.000000,0.185069),(2.000000,2.012875),(3.000000,-0.920342),(3.000000,-0.067366),(3.000000,0.785611),(4.000000,-0.821356),(4.000000,-0.288246),(4.000000,0.244864),(5.000000,-1.235766),(5.000000,0.018611),(5.000000,1.022113),(6.000000,-1.604131),(6.000000,0.054434),(6.000000,1.713000),(7.000000,-1.260318),(7.000000,0.086487),(7.000000,0.984357),(8.000000,-0.737642),(8.000000,0.541823),(8.000000,1.821287),(9.000000,-1.686372),(9.000000,-0.467834),(9.000000,0.953793)}

\newenvironment{seminarvisual}{\par\medskip\noindent\begin{minipage}{\linewidth}\centering}
  {\end{minipage}\par\medskip}
\newcommand{\visualcaption}[1]{\par\smallskip{\small #1\par}}

\newcommand{\olsdistancevisual}{%
\begin{seminarvisual}
\begin{tikzpicture}[x=.9cm,y=.6cm,font=\small]
  % Illustrative data: x=(1,2,3,4,5), y=(2.6,2.4,4,4.4,6.6).
  % The exact OLS fit is yhat=1+x; SSE=1.44. Alternative: 2+0.5*x; SSE=5.19.
  \foreach \shift/\title/\aa/\bb/\col in
    {0/{Другая прямая}/2/.5/accent,7/{Прямая МНК}/1/1/solbar} {
    \begin{scope}[xshift=\shift cm]
      \node[font=\small\bfseries] at (3,8.2) {\title};
      \ifnum\shift=0
        \node at (3,7.55) {$\widetilde y=2+0{,}5x$};
      \else
        \node at (3,7.55) {$\widehat y=1+x$};
      \fi
      \draw[->] (0,0)--(6.25,0) node[right] {$x$};
      \draw[->] (0,0)--(0,7.1) node[above] {$y$};
      \foreach \xx in {1,...,5}
        \draw (\xx,0)--(\xx,-.12) node[below] {\xx};
      \foreach \yy in {0,2,4,6}
        \draw (0,\yy)--(-.08,\yy) node[left] {\yy};
      \draw[\col,very thick] (0,\aa)--(5.6,{\aa+\bb*5.6});
      \foreach \xx/\yy in {1/2.6,2/2.4,3/4,4/4.4,5/6.6} {
        \draw[\col,thick,dashed] (\xx,{\aa+\bb*\xx})--(\xx,\yy);
        \draw[\col,fill=white] (\xx,{\aa+\bb*\xx}) circle (2pt);
        \fill[black] (\xx,\yy) circle (2.5pt);
      }
      \ifnum\shift=0
        \node[anchor=west,text=accent] at (5.12,5.5) {$2{,}1$};
        \node[font=\small\bfseries,text=accent] at (3,-1.5)
          {$\sum_t(y_t-\widetilde y_t)^2=5{,}19$};
      \else
        \node[anchor=west,text=solbar] at (5.25,6.9) {$0{,}6$};
        \node[font=\small\bfseries,text=solbar] at (3,-1.5)
          {$\sum_t(y_t-\widehat y_t)^2=1{,}44\;\text{(минимум)}$};
      \fi
    \end{scope}
  }
\end{tikzpicture}
\visualcaption{Учебный пример: на обоих рисунках одни и те же пять наблюдений
(чёрные точки). Пустые кружки — прогнозы при тех же $x_t$;
пунктирные вертикальные отрезки соединяют наблюдение с прогнозом.}
\end{seminarvisual}}

\newcommand{\variancevisual}{%
\begin{seminarvisual}
\begin{tikzpicture}[x=1cm,y=1cm,font=\small]
  \foreach \shift/\title in {0/{Постоянный разброс},7/{Растущий разброс}} {
    \begin{scope}[xshift=\shift cm]
      \node[font=\small\bfseries] at (2.7,1.75) {\title};
      \draw[->] (0,-1.3)--(0,1.4) node[left] {$e$};
      \draw[->,gray] (0,0)--(5.5,0) node[right] {$x$};
      \node[left] at (0,0) {0};
      \foreach \xx/\rr in {0.35/0.5,0.65/-0.7,0.95/0.2,1.25/0.8,1.55/-0.4,1.85/0.1,2.15/-0.8,2.45/0.6,2.75/-0.2,3.05/0.7,3.35/-0.6,3.65/0.1,3.95/-0.9,4.25/0.8,4.55/0.3,4.85/-0.5,5.15/0.9} {
        \ifnum\shift=0
          \fill[accent] (\xx,\rr) circle (2pt);
        \else
          \fill[accent] (\xx,{\rr*(0.2+0.23*\xx)}) circle (2pt);
        \fi
      }
      \ifnum\shift=0
        \draw[accent,dashed] (0.2,1)--(5.3,1);
        \draw[accent,dashed] (0.2,-1)--(5.3,-1);
      \else
        \draw[accent,dashed] (0.2,0.246)--(5.3,1.419);
        \draw[accent,dashed] (0.2,-0.246)--(5.3,-1.419);
      \fi
    \end{scope}
  }
\end{tikzpicture}
\visualcaption{Схематический пример: важна ширина облака вокруг нуля,
а не наклон линии регрессии. Пунктир показывает изменение масштаба;
это не доверительные границы.}
\end{seminarvisual}}

\newcommand{\fglsvisual}{%
\begin{seminarvisual}
\begin{tikzpicture}[font=\small,box/.style={draw=accent,rounded corners=2pt,
  fill=accent!5,align=center,text width=3.5cm,minimum height=1.15cm}]
  \node[box] (a) at (0,0) {МНК исходной модели\\$e_t=y_t-\widehat y_t$};
  \node[box] (b) at (4.6,0) {Регрессия $e_t^2$ на $1,x_t$\\$\widehat\sigma_t^2=\widehat\mu+\widehat\delta x_t$};
  \node[box,draw=solbar,fill=solbar!6] (c) at (9.2,0) {МНК после деления\\на $\widehat\sigma_t$};
  \draw[->,thick,accent] (a)--(b);
  \draw[->,thick,solbar] (b)--(c);
  \draw[accent] (-1.85,0.85)--(-1.85,1.02)--(6.45,1.02)--(6.45,0.85);
  \node[above,font=\small\bfseries] at (2.3,1.02) {Этап 1: оценить дисперсии};
  \node[above,font=\small\bfseries] at (9.2,1.02) {Этап 2: ВМНК};
\end{tikzpicture}
\visualcaption{Два этапа включают три регрессии. Во втором этапе вес
наблюдения равен $1/\widehat\sigma_t^2$; делитель — $\widehat\sigma_t$.}
\end{seminarvisual}}

\newcommand{\dwvisual}{%
\begin{seminarvisual}
\begin{tikzpicture}[x=3cm,y=1cm,font=\small]
  \fill[accent!15] (0,0) rectangle (1.27,.6);
  \fill[notebar!25] (1.27,0) rectangle (1.45,.6);
  \fill[solbar!12] (1.45,0) rectangle (4,.6);
  \draw[->] (0,0)--(4.1,0) node[right] {DW};
  \draw[dashed] (1.27,0)--(1.27,.8);
  \draw[dashed] (1.45,0)--(1.45,.8);
  \node[anchor=north east] at (1.27,-.12) {$d_L=1{,}27$};
  \node[anchor=north west] at (1.45,-.12) {$d_U=1{,}45$};
  \node[below] at (0,0) {0}; \node[below] at (4,0) {4};
  \node at (.60,.3) {Отвергаем $H_0$};
  \node at (2.75,.3) {Не отвергаем $H_0$};
  \node[anchor=south west] at (1.5,1.25) {Зона неопределённости};
  \draw[->,notebar] (1.65,1.2)--(1.36,.65);
  \draw[->,very thick,accent] (1.103,1.08)--(1.103,.64);
  \node[anchor=south east,text=accent] at (1.13,1.08) {$\mathrm{DW}=1{,}103$};
\end{tikzpicture}
\visualcaption{Задача 6.6: $H_0\!:\rho=0$ против $H_1\!:\rho>0$,
уровень 5\,\%. Наблюдаемая статистика попадает в область отвержения.
Правая область относится только к этой односторонней проверке.}
\end{seminarvisual}}

\newcommand{\arresidualvisual}{%
\begin{seminarvisual}
\begin{tikzpicture}[x=.235cm,y=2.3cm,font=\small]
  \foreach \shift/\title in {0/{Исходные остатки $e_t$},30/{Инновации после AR(1)}} {
    \begin{scope}[xshift=\shift*.235cm]
      \node[font=\small\bfseries] at (12.5,.88) {\title};
      \draw[->] (0,-.65)--(0,.70);
      \draw[gray,dashed] (0,0)--(24,0);
      \draw[->] (0,-.65)--(25,-.65) node[right] {$t$};
      \foreach \xx in {1,6,12,18,24}
        \draw (\xx,-.65)--(\xx,-.69) node[below] {\xx};
      \foreach \yy/\lab in {-.5/{-0{,}5},0/0,.5/{0{,}5}}
        \node[left] at (0,\yy) {$\lab$};
      \ifnum\shift=0
        \draw[accent,thick] plot[mark=*,mark size=1.7pt] coordinates {\olsTimeCoordinates};
      \else
        \draw[solbar,thick] plot[mark=*,mark size=1.7pt] coordinates {\arTimeCoordinates};
      \fi
    \end{scope}
  }
\end{tikzpicture}
\visualcaption{Данные задачи 6.6, одинаковый масштаб по вертикали.
Слева $\mathrm{DW}=1{,}103$, справа $2{,}162$.
Справа показаны $\widehat\eta_t=\widehat\varepsilon_t-
\widehat\rho\widehat\varepsilon_{t-1}$, поэтому ряд начинается с $t=2$.
График дополняет диагностику, но не доказывает независимость.}
\end{seminarvisual}}

\newcommand{\intervalvisual}{%
\begin{seminarvisual}
\begin{tikzpicture}[x=14cm,y=1cm,font=\small]
  % Plot x = beta + 1.1 to keep node positions compact.
  \draw[->] (0,0)--(.63,0) node[right] {$\beta_2$};
  \foreach \xx/\lab in {0/{-1{,}1},.1/{-1{,}0},.2/{-0{,}9},.3/{-0{,}8},.4/{-0{,}7},.5/{-0{,}6},.6/{-0{,}5}}
    \draw (\xx,0)--(\xx,-.07) node[below] {$\lab$};
  \node[anchor=east,align=right] at (-.02,1.65) {МНК};
  \node[anchor=east,align=right] at (-.02,.65) {С AR(1)};
  \draw[accent,line width=1.3pt] (.086,1.65)--(.556,1.65);
  \draw[solbar,line width=1.3pt] (.075,.65)--(.462,.65);
  \foreach \xx in {.086,.556} \draw[accent] (\xx,1.55)--(\xx,1.75);
  \foreach \xx in {.075,.462} \draw[solbar] (\xx,.55)--(\xx,.75);
  \fill[accent] (.3209,1.65) circle (2.7pt);
  \fill[solbar] (.2688,.65) circle (2.7pt);
  \node[above] at (.086,1.75) {$-1{,}014$};
  \node[above] at (.556,1.75) {$-0{,}544$};
  \node[above] at (.075,.75) {$-1{,}025$};
  \node[above] at (.462,.75) {$-0{,}638$};
\end{tikzpicture}
\visualcaption{Точки — оценки $\beta_2$, отрезки — 95\,\%-ные интервалы.
Интервал МНК построен при нарушенной предпосылке независимости;
интервал с AR(1) — приближённый.}
\end{seminarvisual}}

\newcommand{\gqvisual}{%
\begin{seminarvisual}
\begin{tikzpicture}[x=.4cm,y=1cm,font=\small]
  \fill[accent!14] (0,0) rectangle (12,.65);
  \fill[gray!15] (12,0) rectangle (18,.65);
  \fill[solbar!15] (18,0) rectangle (30,.65);
  \foreach \xx in {0,...,30} \draw[white,thick] (\xx,0)--(\xx,.65);
  \node[above,align=center] at (6,.7) {Первые 12 семей\\$Y=12000$--$15000$};
  \node[above,align=center] at (15,.7) {Пропускаем 6\\$16000$--$17000$};
  \node[above,align=center] at (24,.7) {Последние 12 семей\\$Y=18000$--$21000$};
  \node[below] at (6,0) {$\ESS_1=1\,046\,000$};
  \node[below] at (24,0) {$\ESS_2=3\,344\,000$};
  \draw[->] (0,-.65)--(30,-.65) node[below,anchor=north east] {Доход возрастает};
\end{tikzpicture}
\visualcaption{Каждая ячейка — одна семья. Оцениваем две отдельные
регрессии $C$ на $Y$ с константой; сравниваем $\ESS_2/10$ и $\ESS_1/10$.}
\end{seminarvisual}}

\newcommand{\wlsresidualvisual}{%
\begin{seminarvisual}
\begin{tikzpicture}[x=.58cm,y=.55cm,font=\small]
  \foreach \shift/\title in {0/{До коррекции},12/{После деления на доход}} {
    \begin{scope}[xshift=\shift*.58cm]
      \node[font=\small\bfseries] at (4.5,3.15) {\title};
      \draw[->] (-.5,-2.6)--(-.5,2.7);
      \draw[gray,dashed] (-.5,0)--(9.5,0);
      \draw[->] (-.5,-2.6)--(9.8,-2.6);
      \foreach \xx/\lab in {0/12,3/15,6/18,9/21}
        \draw (\xx,-2.6)--(\xx,-2.72) node[below] {\lab};
      \foreach \yy in {-2,0,2} \node[left] at (-.5,\yy) {$\yy$};
      \node[below] at (4.5,-3.3) {Доход, тыс. долл.};
      \ifnum\shift=0
        \foreach \p in \olsIncomeCoordinates \fill[accent] \p circle (2pt);
      \else
        \foreach \p in \wlsIncomeCoordinates \fill[solbar] \p circle (2pt);
      \fi
    \end{scope}
  }
\end{tikzpicture}
\visualcaption{Данные задачи 6.7: слева остатки МНК $e_t$, справа остатки
ВМНК $\widehat v_t=(C_t-\widehat\alpha-\widehat\beta Y_t)/Y_t$. Чтобы
ряды были в одном масштабе, каждый разделён на корень из среднего квадрата
своих значений. Слева у 15 семей с высоким доходом разброс почти вдвое
больше, чем у 15 семей с низким (1,24 против 0,69); справа разница
заметно меньше (1,12 против 0,87).}
\end{seminarvisual}}

\newcommand{\spearmanrankvisual}{%
\begin{seminarvisual}
\begin{tikzpicture}[x=.93cm,y=.52cm,font=\scriptsize]
  % Problem 6.7: |e_t| sorted ascending; two strips of 15 positions.
  \foreach \p/\v/\r/\fam/\tie in {0/$8{,}48$/$1$/17/0,1/$25{,}45$/$2{,}5$/11/1,2/$25{,}45$/$2{,}5$/20/1,3/$42{,}42$/$4$/23/0,4/$57{,}58$/$5$/8/0,5/$59{,}39$/$6$/5/0,6/$76{,}36$/$7$/2/0,7/$91{,}52$/$8$/15/0,8/$108{,}48$/$9$/14/0,9/$223{,}64$/$10{,}5$/3/1,10/$223{,}64$/$10{,}5$/29/1,11/$242{,}42$/$12$/7/0,12/$259{,}39$/$13$/26/0,13/$274{,}55$/$14$/12/0,14/$276{,}36$/$15$/1/0,15/$308{,}48$/$16$/13/0,16/$325{,}45$/$17$/10/0,17/$340{,}61$/$18{,}5$/6/1,18/$340{,}61$/$18{,}5$/25/1,19/$359{,}39$/$20$/4/0,20/$408{,}48$/$21$/18/0,21/$442{,}42$/$22$/24/0,22/$476{,}36$/$23$/30/0,23/$491{,}52$/$24$/16/0,24/$557{,}58$/$25$/22/0,25/$657{,}58$/$26$/9/0,26/$674{,}55$/$27$/19/0,27/$725{,}45$/$28$/21/0,28/$823{,}64$/$29$/28/0,29/$859{,}39$/$30$/27/0} {
    \pgfmathtruncatemacro{\row}{\p>14 ? 1 : 0}
    \pgfmathsetmacro{\xx}{\p-15*\row}
    \pgfmathsetmacro{\yy}{-4.2*\row}
    \begin{scope}[shift={(\xx,\yy)}]
      \ifnum\tie=1 \fill[notebar!35] (0,1) rectangle (1,3);
      \else \fill[gray!10] (0,1) rectangle (1,3); \fi
      \node at (.5,2.5) {\v};
      \node at (.5,1.5) {\r};
      \node[text=gray] at (.5,.5) {\fam};
      \draw[white,thick] (0,1)--(0,3) (1,1)--(1,3);
    \end{scope}
  }
  \foreach \yy in {0,-4.2} {
    \node[anchor=east,font=\small] at (-.05,2.5+\yy) {$|e_t|$};
    \node[anchor=east,font=\small] at (-.05,1.5+\yy) {ранг};
    \node[anchor=east,font=\small,text=gray] at (-.05,.5+\yy) {семья};
  }
  \draw[->] (0,-4.55)--(15,-4.55) node[below,anchor=north east,font=\small] {$|e_t|$ возрастает};
\end{tikzpicture}
\visualcaption{30 значений $|e_t|$ по возрастанию (вторая строка рисунка —
продолжение первой). Ранг — номер места; выделенные цветом ячейки — совпадающие
значения, им дан средний ранг занимаемых мест.}
\end{seminarvisual}}

\newcommand{\spearmanvisual}{%
\begin{seminarvisual}
\begin{tikzpicture}[x=.93cm,y=.52cm,font=\scriptsize]
  % Problem 6.7: families t=1..30 in income order; two strips of 15.
  \foreach \k in {0,...,9} {
    \pgfmathtruncatemacro{\inc}{12+\k}
    \pgfmathtruncatemacro{\ry}{3*\k+2}
    \pgfmathtruncatemacro{\row}{\k>4 ? 1 : 0}
    \pgfmathsetmacro{\xx}{3*\k-15*\row}
    \pgfmathsetmacro{\yy}{-4.8*\row}
    \ifodd\k \def\blk{gray!22}\else\def\blk{gray!10}\fi
    \fill[\blk] (\xx,2+\yy) rectangle (\xx+3,3.4+\yy);
    \node at (\xx+1.5,3.0+\yy) {$Y=\inc\,000$};
    \node at (\xx+1.5,2.4+\yy) {$R_Y=\ry$};
    \draw[white,thick] (\xx,2+\yy)--(\xx,3.4+\yy);
  }
  \foreach \i/\r/\s in {0/$15$/36,1/$7$/20,2/$10{,}5$/27,3/$20$/46,4/$6$/18,5/$18{,}5$/43,6/$12$/30,7/$5$/16,8/$26$/58,9/$17$/40,10/$2{,}5$/11,11/$14$/34,12/$16$/38,13/$9$/24,14/$8$/22,15/$24$/54,16/$1$/8,17/$21$/48,18/$27$/60,19/$2{,}5$/11,20/$28$/62,21/$25$/56,22/$4$/14,23/$22$/50,24/$18{,}5$/43,25/$13$/32,26/$30$/66,27/$29$/64,28/$10{,}5$/27,29/$23$/52} {
    \pgfmathtruncatemacro{\row}{\i>14 ? 1 : 0}
    \pgfmathsetmacro{\xx}{\i-15*\row}
    \pgfmathsetmacro{\yy}{-4.8*\row}
    \fill[accent!\s] (\xx,1+\yy) rectangle (\xx+1,2+\yy);
    \node at (\xx+.5,1.5+\yy) {\r};
    \draw[white,thick] (\xx,1+\yy)--(\xx,2+\yy);
  }
  \foreach \i/\dd in {0/$-13$,1/$-5$,2/$-8{,}5$,3/$-15$,4/$-1$,5/$-13{,}5$,6/$-4$,7/$3$,8/$-18$,9/$-6$,10/$8{,}5$,11/$-3$,12/$-2$,13/$5$,14/$6$,15/$-7$,16/$16$,17/$-4$,18/$-7$,19/$17{,}5$,20/$-8$,21/$-2$,22/$19$,23/$1$,24/$7{,}5$,25/$13$,26/$-4$,27/$0$,28/$18{,}5$,29/$6$} {
    \pgfmathtruncatemacro{\row}{\i>14 ? 1 : 0}
    \pgfmathsetmacro{\xx}{\i-15*\row}
    \pgfmathsetmacro{\yy}{-4.8*\row}
    \node at (\xx+.5,.5+\yy) {\dd};
    \node[text=gray] at (\xx+.5,-.4+\yy) {\pgfmathparse{int(\i+1)}\pgfmathresult};
  }
  \foreach \yy in {0,-4.8} {
    \node[anchor=east,font=\small] at (-.05,2.7+\yy) {доход};
    \node[anchor=east,font=\small] at (-.05,1.5+\yy) {$R_{|e|}$};
    \node[anchor=east,font=\small] at (-.05,.5+\yy) {$d_t$};
    \node[anchor=east,font=\small,text=gray] at (-.05,-.4+\yy) {семья};
  }
\end{tikzpicture}
\visualcaption{Семьи в порядке роста дохода (вторая строка — продолжение
первой). Чем больше ранг $|e_t|$, тем темнее ячейка; $d_t=R_Y-R_{|e|}$,
$\sum d_t^2=2931{,}5$. Тёмные ячейки чаще во второй строке, но не только
там: связь положительная, но слабая.}
\end{seminarvisual}}

\section*{Напоминание: МНК и его предпосылки}

\textbf{Модель и МНК.} Парная регрессия
\[
  y_t = \alpha + \beta x_t + \varepsilon_t, \qquad t = 1, \dots, n,
\]
где $y_t$ — зависимая переменная, $x_t$ — регрессор, $\alpha$ и $\beta$ —
неизвестные параметры, $\varepsilon_t$ — ненаблюдаемая случайная ошибка,
$n$ — число наблюдений. МНК выбирает оценки $\widehat\alpha$, $\widehat\beta$,
при которых сумма квадратов отклонений наблюдений от прямой наименьшая:
\[
  (\widehat\alpha,\widehat\beta) = \arg\min_{a,\,b}\sum_{t=1}^n(y_t - a - bx_t)^2,
  \qquad
  e_t = y_t - \widehat\alpha - \widehat\beta x_t ,
\]
где $a$, $b$ — пробные значения свободного члена и наклона, $e_t$ —
\emph{остаток}: отклонение наблюдения от оценённой прямой. Остаток вычисляется
по данным и служит оценкой ненаблюдаемой ошибки $\varepsilon_t$.

\olsdistancevisual

Отклонения измеряются по вертикали, при том же $x_t$. Например, справа для
пятой точки $e_5 = 6{,}6 - 6 = 0{,}6$, и её вклад в сумму равен $0{,}36$.
У прямой МНК сумма квадратов меньше, чем у любой другой прямой
($1{,}44$ против $5{,}19$ слева).

\textbf{Суммы квадратов.} $\TSS = \sum(y_t - \overline{y})^2$ — полная,
$\RSS = \sum(\widehat y_t - \overline{y})^2$ — объяснённая регрессией,
$\ESS = \sum e_t^2$ — остаточная сумма квадратов; $\overline{y}$ — среднее
$y_t$, $\widehat y_t = \widehat\alpha + \widehat\beta x_t$ — подогнанные значения.
$\TSS = \RSS + \ESS$, коэффициент детерминации $R^2 = \RSS/\TSS$.

\textbf{Обозначения.} $\E$ — математическое ожидание, $\D$ — дисперсия,
$\Cov$ — ковариация; $k$ — число оцениваемых параметров вместе с константой
(в парной регрессии $k = 2$); $F_\alpha$, $\chi^2_\alpha$, $t_\alpha$ —
критические значения распределений для уровня значимости $\alpha$.

\textbf{Предпосылки} (значения регрессоров считаем фиксированными):
\begin{enumerate}[label=\arabic*),leftmargin=2.2em,itemsep=2pt]
  \item модель линейна по параметрам, и не все $x_t$ одинаковы (в
        множественной регрессии — ни один регрессор не является точной
        линейной комбинацией остальных);
  \item $\E(\varepsilon_t) = 0$ — в среднем ошибка равна нулю;
  \item $\D(\varepsilon_t) = \sigma^2$ для всех $t$ (\emph{гомоскедастичность})
        и $\Cov(\varepsilon_t, \varepsilon_s) = 0$ при $t \ne s$ (нет
        \emph{автокорреляции}); $\sigma^2$ — общая для всех наблюдений
        дисперсия ошибки;
  \item ошибки нормально распределены — нужна только для точных $t$- и
        $F$-тестов в малых выборках.
\end{enumerate}

\begin{center}
\small
\begin{tabular}{@{}>{\raggedright\arraybackslash}p{4.4cm}>{\raggedright\arraybackslash}p{2.6cm}>{\raggedright\arraybackslash}p{7.6cm}@{}}
\toprule
Свойство МНК & Предпосылки & Смысл \\
\midrule
Несмещённость & 1--2 & В среднем по многим выборкам оценки равны истинным параметрам. \\
Эффективность (BLUE) & 1--3 & Наименьшая дисперсия среди линейных несмещённых оценок (теорема Гаусса--Маркова). \\
Верные ст.~ошибки, $t$- и $F$-тесты & 1--3; точные — 1--4 & Интервалы и тесты имеют заявленный уровень. \\
\bottomrule
\end{tabular}
\end{center}

\textbf{Что в этом семинаре.} Нарушается предпосылка 3: гетероскедастичность
— её первая часть, автокорреляция — вторая. Предпосылки 1--2 сохраняются,
поэтому МНК-оценки остаются несмещёнными, но перестают быть эффективными, а
обычные стандартные ошибки, $t$- и $F$-тесты и интервалы становятся неверными.
План: обнаружить нарушение по остаткам, затем преобразовать модель так, чтобы
предпосылка 3 снова выполнялась.

\clearpage
\section{Основные формулы}

% -------------------------------------------------------------------
\subsection{Гетероскедастичность}\label{sec:het}

\emph{Гетероскедастичность} — дисперсия ошибки различна в разных
наблюдениях: $\D(\varepsilon_t) = \sigma_t^2$, где $\sigma_t^2$ — дисперсия
ошибки в наблюдении $t$, а $\sigma_t$ — её стандартное отклонение. Обычно
разброс растёт с масштабом: например, расходы богатых семей колеблются
сильнее, чем бедных.

\variancevisual

\textbf{Как обнаружить.} Ошибки не наблюдаются, поэтому все методы работают с
остатками МНК: $e_t^2$ и $|e_t|$ — грубые оценки дисперсии $\sigma_t^2$ и
стандартного отклонения $\sigma_t$. Во всех тестах $H_0$ —
гомоскедастичность (все $\sigma_t^2$ равны).
\begin{enumerate}[label=\arabic*), leftmargin=2.2em, itemsep=4pt]
  \item \emph{График} $e_t^2$ (или $|e_t|$) против регрессора $x_t$.
        Расширяющийся «веер» — признак гетероскедастичности. Это подсказка,
        а не тест.

  \item \emph{Тест Голдфелда--Квандта} — когда дисперсия, предположительно,
        растёт с $x_t$. Упорядочиваем наблюдения по возрастанию $x_t$,
        выбрасываем $c$ средних и оцениваем регрессию отдельно по первым $m$
        и по последним $m$ наблюдениям, $m = (n - c)/2$:
        \begin{equation}
          F = \frac{\ESS_2/(m-k)}{\ESS_1/(m-k)} = \frac{\ESS_2}{\ESS_1}
          \sim F(m-k,\ m-k),
          \label{eq:gq}
        \end{equation}
        где $\ESS_1$ и $\ESS_2$ — остаточные суммы квадратов регрессий по
        группе с малыми и с большими $x_t$, $m - k$ — число степеней свободы
        в каждой. Альтернатива $H_1$: в группе с большими $x_t$ дисперсия
        больше. $H_0$ отвергаем, если $F > F_\alpha(m-k,\ m-k)$. Средние
        наблюдения выбрасываем, чтобы группы сильнее различались.

  \item \emph{Тест Бройша--Пагана.} Предполагаем, что дисперсия может
        линейно зависеть от наблюдаемых переменных $z_{1t}, \dots, z_{pt}$
        (обычно это регрессоры модели):
        \[
          \sigma_t^2 = \sigma^2\bigl(\gamma_0 + \gamma_1z_{1t} + \dots + \gamma_pz_{pt}\bigr),
          \qquad
          H_0\colon \gamma_1 = \dots = \gamma_p = 0,
        \]
        где $\gamma_j$ — неизвестные коэффициенты, $p$ — число переменных $z$.
        При $H_0$ дисперсия постоянна. Так как $\E\varepsilon_t^2 = \sigma_t^2$,
        проверяем, объясняют ли $z$ разброс квадратов остатков. Расчёт:
        \begin{enumerate}[label=\asbuk*), leftmargin=1.8em, itemsep=1pt]
          \item оцениваем исходную модель МНК, получаем остатки $e_t$;
          \item $\widehat{\sigma}^2 = \ESS/n$ — средний квадрат остатка
                (делим именно на $n$);
          \item $g_t = e_t^2/\widehat{\sigma}^2$ — нормированные квадраты
                остатков, их среднее $\overline{g} = 1$;
          \item оцениваем вспомогательную регрессию
                $g_t = a_0 + a_1z_{1t} + \dots + a_pz_{pt} + v_t$ ($v_t$ — её
                ошибка) и берём её \emph{объяснённую} сумму квадратов
                $\RSS_{\text{aux}} = \sum(\widehat g_t - \overline{g})^2$, где
                $\widehat g_t$ — подогнанные значения;
          \item вычисляем
                \begin{equation}
                  \mathrm{BP} = \frac{\RSS_{\text{aux}}}{2} \overset{a}{\sim} \chi^2(p)
                  \label{eq:bp}
                \end{equation}
                и отвергаем $H_0$, если $\mathrm{BP} > \chi^2_\alpha(p)$.
        \end{enumerate}
        Знак $\overset{a}{\sim}$ означает, что распределение приближённое
        (точное лишь при $n \to \infty$). Деление на 2 нужно потому, что при
        $H_0$ и нормальных ошибках $\varepsilon_t^2/\sigma^2 \sim \chi^2(1)$,
        а у $\chi^2(1)$ дисперсия равна 2: без деления статистика была бы
        вдвое больше, чем $\chi^2(p)$.

  \item \emph{Тест Уайта.} Для модели с двумя регрессорами
        $y_t = \beta_1 + \beta_2x_{2t} + \beta_3x_{3t} + \varepsilon_t$
        оцениваем вспомогательную регрессию квадратов остатков на регрессоры,
        их квадраты и произведение:
        \[
          e_t^2 = a_0 + a_1x_{2t} + a_2x_{3t} + a_3x_{2t}^2 + a_4x_{3t}^2
                + a_5x_{2t}x_{3t} + v_t ,
          \qquad
          H_0\colon a_1 = \dots = a_5 = 0,
        \]
        и вычисляем
        \begin{equation}
          W = nR^2_{\text{aux}} \overset{a}{\sim} \chi^2(p),
          \qquad
          R^2_{\text{aux}} = 1 - \frac{\sum_t\widehat v_t^{\,2}}
                                      {\sum_t\bigl(e_t^2 - \overline{e^2}\bigr)^2},
          \label{eq:white}
        \end{equation}
        где $R^2_{\text{aux}}$ — коэффициент детерминации
        \emph{вспомогательной} регрессии (доля разброса $e_t^2$, объяснённая
        её регрессорами), $\widehat v_t$ — её остатки, $\overline{e^2}$ —
        среднее $e_t^2$, $p$ — число её регрессоров без константы (здесь
        $p = 5$; в парной регрессии остаются $x_t$ и $x_t^2$, $p = 2$).
        $H_0$ отвергаем, если $W > \chi^2_\alpha(p)$. Квадраты и произведения
        приближают любую гладкую зависимость $\sigma_t^2$ от регрессоров,
        поэтому форму гетероскедастичности знать не нужно; цена — много
        степеней свободы и меньшая мощность при малом $n$. Если вместо них
        взять те же $z_t$, что в тесте Бройша--Пагана, получится его вариант
        Кёнкера $nR^2_{\text{aux}} \overset{a}{\sim} \chi^2(p)$, не требующий
        нормальности ошибок.

  \item \emph{Тест Глейзера} — ищет форму зависимости $\sigma_t$ от $x_t$.
        Для нескольких значений показателя $\gamma$ (например, $1/2$, $1$, $2$)
        оцениваем регрессию
        \[
          |e_t| = a_0 + a_1x_t^{\gamma} + v_t ,
          \qquad
          H_0\colon a_1 = 0,
          \qquad
          t = \frac{\widehat a_1}{s_{\widehat a_1}} \overset{a}{\sim} t(n-2),
        \]
        где $\widehat a_1$ — оценка наклона, $s_{\widehat a_1}$ — её
        стандартная ошибка. Значимый положительный наклон означает, что
        разброс ошибок растёт с $x_t$. Лучше всего подогнанная форма
        подсказывает веса для взвешенного МНК: например, если лучше всех
        подходит $\gamma = 1$, то $\sigma_t \propto x_t$ и веса равны $1/x_t$.

  \item \emph{Ранговая корреляция Спирмена} между $|e_t|$ и $x_t$:
        \begin{equation}
          r_S = 1 - \frac{6\sum d_t^2}{n(n^2 - 1)},
          \qquad
          t = \frac{r_S\sqrt{n-2}}{\sqrt{1 - r_S^2}} \overset{a}{\sim} t(n-2),
          \label{eq:spearman}
        \end{equation}
        где $d_t$ — разность рангов $|e_t|$ и $x_t$ в наблюдении $t$ (ранг —
        номер значения в ряду, упорядоченном по возрастанию). При совпадающих
        значениях им дают средний ранг, и $r_S$ считают как обычный
        коэффициент корреляции двух столбцов рангов. Значимая положительная
        $r_S$ означает, что разброс растёт с $x_t$.
\end{enumerate}

\textbf{Коррекция: взвешенный МНК.} Пусть дисперсия известна с точностью до
множителя: $\sigma_t^2 = \sigma^2z_t^2$, где $z_t > 0$ — известная величина
(например, $z_t = x_t$). Делим уравнение на $z_t$:
\begin{equation}
  \frac{y_t}{z_t} = \alpha\frac{1}{z_t} + \beta\frac{x_t}{z_t} + \frac{\varepsilon_t}{z_t},
  \qquad
  \D\Bigl(\frac{\varepsilon_t}{z_t}\Bigr) = \frac{\D(\varepsilon_t)}{z_t^2} = \sigma^2 .
  \label{eq:wls}
\end{equation}
Ошибка нового уравнения гомоскедастична, поэтому его оценивают обычным МНК.
Свободного члена в нём нет: $\alpha$ стоит при $1/z_t$. При $z_t = x_t$
получается $y_t/x_t = \alpha/x_t + \beta + \varepsilon_t/x_t$, и уже $\beta$
играет роль свободного члена. Деление на $z_t$ означает, что в сумме квадратов
наблюдение получает вес $1/z_t^2$: чем больше разброс, тем меньше вес. Если
$z_t$ неизвестно, его сначала оценивают по остаткам (\emph{доступный}
обобщённый МНК, задача 6.5).

% -------------------------------------------------------------------
\subsection{Автокорреляция ошибок}\label{sec:ac}

\emph{Автокорреляция} — ошибки разных наблюдений коррелированы; обычно
встречается во временных рядах, где $t$ — время. Простейшая модель —
автокорреляция первого порядка AR(1):
\[
  \varepsilon_t = \rho\varepsilon_{t-1} + \eta_t, \qquad |\rho| < 1,
\]
где $\rho$ — коэффициент автокорреляции (корреляция соседних ошибок), а
$\eta_t$ — независимые случайные величины $N(0, \sigma_\eta^2)$, «новая» часть
ошибки в момент $t$ (её называют инновацией). При $\rho > 0$ соседние ошибки
обычно одного знака, при $\rho = 0$ автокорреляции нет. Последствия те же, что
при гетероскедастичности: оценки МНК несмещённые, но неэффективные, а обычные
стандартные ошибки неверны.

\textbf{Статистика Дарбина--Уотсона:}
\begin{equation}
  \mathrm{DW} = \frac{\sum_{t=2}^{n}(e_t - e_{t-1})^2}{\sum_{t=1}^{n}e_t^2}
  \approx 2(1 - \widehat{\rho}),
  \qquad
  \widehat{\rho} = \frac{\sum_{t=2}^{n}e_te_{t-1}}{\sum_{t=2}^{n}e_{t-1}^2},
  \label{eq:dw}
\end{equation}
где $e_t$ — остатки МНК, взятые в исходном временном порядке, $\widehat\rho$
— оценка $\rho$ по ним. Если соседние остатки похожи ($\rho > 0$), разности
$e_t - e_{t-1}$ малы и DW близка к 0; при $\rho = 0$ $\mathrm{DW} \approx 2$;
при $\rho < 0$ DW близка к 4. Гипотезу $H_0\colon \rho = 0$ против
$H_1\colon \rho > 0$ проверяют по табличным границам $d_L < d_U$, которые
зависят от $n$ и числа регрессоров:
\[
\begin{array}{c|c|c}
\mathrm{DW}<d_L & d_L\leq\mathrm{DW}\leq d_U & \mathrm{DW}>d_U\\
\text{отвергаем }H_0 & \text{вывода нет} & \text{не отвергаем }H_0
\end{array}
\]
Таблицы DW применимы к регрессии с константой, в которой среди регрессоров нет
лагированной зависимой переменной $y_{t-1}$.

\textbf{Коррекция: процедура Кохрейна--Оркатта.} Если $\rho$ известно,
вычтем из уравнения $y_t = \beta_1 + \beta_2x_t + \varepsilon_t$ то же
уравнение для момента $t-1$, умноженное на $\rho$:
\begin{equation}
  y_t - \rho y_{t-1} = \beta_1(1 - \rho) + \beta_2(x_t - \rho x_{t-1}) + \eta_t,
  \qquad t = 2, \dots, n .
  \label{eq:co}
\end{equation}
Величины $y_t - \rho y_{t-1}$ и $x_t - \rho x_{t-1}$ называют
\emph{квазиразностями}. Ошибка $\eta_t$ уже не автокоррелирована, поэтому
уравнение оценивают обычным МНК с константой; первое наблюдение теряется.
Константа этого уравнения равна $\beta_1(1 - \rho)$, поэтому
$\widehat\beta_1$ получают делением её оценки на $1 - \rho$, а наклон сразу
даёт $\widehat\beta_2$. На практике $\rho$ неизвестно, и шаги повторяют:
1) МНК исходного уравнения и остатки $e_t$; 2) $\widehat\rho$ по \eqref{eq:dw};
3) МНК для квазиразностей с $\widehat\rho$, оценки $\widehat\beta_1$,
$\widehat\beta_2$; 4) новые остатки \emph{исходного} уравнения
$e_t = y_t - \widehat\beta_1 - \widehat\beta_2x_t$ и возврат к шагу 2 — пока
$\widehat\rho$ не перестанет меняться.
\fi

% ===================================================================
\ifsolutions\clearpage\fi
\section{Задачи}

% -------------------------------------------------------------------
\problem{6.5}

Предположим, что для системы $y_t = \alpha + \beta x_t + \varepsilon_t$,
$t = 1, \dots, n$, выполнены все предположения классической нормальной
модели, за одним исключением: дисперсии ошибок удовлетворяют соотношениям
\[
  \sigma_t^2 = \mu + \delta x_t .
\]
Предложите двухшаговую процедуру оценивания параметров $\alpha$ и $\beta$.

\begin{solution}
\textbf{Что не так и что хотим получить.} Дисперсия ошибки меняется с $x_t$,
поэтому обычный МНК для исходного уравнения даёт несмещённые, но
неэффективные оценки, а его стандартные ошибки неверны. Если бы мы знали
стандартные отклонения ошибок $\sigma_t$, то разделили бы уравнение на
$\sigma_t$ по \eqref{eq:wls} (с $z_t = \sigma_t$): у новой ошибки
$\varepsilon_t/\sigma_t$ дисперсия равна $\sigma_t^2/\sigma_t^2 = 1$ для всех
$t$, и обычный МНК снова оптимален. Но $\sigma_t$ неизвестны. Зато известна
\emph{форма} дисперсии: $\sigma_t^2 = \mu + \delta x_t$ — она зависит всего от
двух неизвестных чисел $\mu$ и $\delta$. Отсюда план из двух этапов:
\textbf{1)} оценить $\mu$ и $\delta$ и по ним — каждое $\sigma_t$;
\textbf{2)} разделить уравнение на оценённые $\widehat\sigma_t$ и применить МНК.

\fglsvisual

\medskip
\textbf{Этап 1. Оцениваем дисперсии ошибок.}

\smallskip
\emph{а) Обычный МНК исходного уравнения.} Оцениваем
$y_t = \alpha + \beta x_t + \varepsilon_t$ и считаем остатки
\[
  e_t = y_t - \widetilde\alpha - \widetilde\beta x_t ,
\]
где $\widetilde\alpha$, $\widetilde\beta$ — оценки обычного МНК (тильда
отличает их от итоговых оценок этапа 2). Эти оценки несмещённые и при
гетероскедастичности, поэтому остатки $e_t$ близки к ненаблюдаемым ошибкам
$\varepsilon_t$.

\smallskip
\emph{б) Почему квадраты остатков можно регрессировать на $x_t$.} Так как
$\E\varepsilon_t = 0$, математическое ожидание квадрата ошибки равно её
дисперсии: $\E\varepsilon_t^2 = \D(\varepsilon_t) = \sigma_t^2 = \mu + \delta x_t$.
Значит, квадрат ошибки можно записать как
\[
  \varepsilon_t^2 = \mu + \delta x_t + v_t ,
  \qquad
  v_t = \varepsilon_t^2 - \sigma_t^2 ,\quad \E v_t = 0 ,
\]
где $v_t$ — случайное отклонение квадрата ошибки от его среднего. Это обычная
линейная регрессия: зависимая переменная $\varepsilon_t^2$, регрессор $x_t$,
параметры $\mu$ и $\delta$. Ошибки $\varepsilon_t$ мы не видим, поэтому
подставляем вместо них остатки из п.~а).

\smallskip
\emph{в) Вспомогательная регрессия.} Оцениваем МНК
\[
  e_t^2 = \mu + \delta x_t + v_t
\]
и получаем $\widehat\mu$, $\widehat\delta$. Отдельный $e_t^2$ — очень шумная
оценка своей дисперсии, но регрессия по всем $n$ наблюдениям усредняет этот
шум и оценивает общую закономерность $\mu + \delta x_t$.

\smallskip
\emph{г) Оценки дисперсий и стандартных отклонений.} Для каждого наблюдения
подставляем $x_t$ в оценённую формулу:
\[
  \widehat\sigma_t^2 = \widehat\mu + \widehat\delta x_t ,
  \qquad
  \widehat\sigma_t = \sqrt{\widehat\mu + \widehat\delta x_t} .
\]
Перед этапом 2 проверяем, что все $\widehat\sigma_t^2 > 0$: иначе корень
не извлечь и делить не на что.

\medskip
\textbf{Этап 2. Взвешенный МНК.} Делим исходное уравнение на
$\widehat\sigma_t$:
\[
  \underbrace{\frac{y_t}{\widehat\sigma_t}}_{y_t^*}
  = \alpha\underbrace{\frac{1}{\widehat\sigma_t}}_{w_{1t}}
  + \beta\underbrace{\frac{x_t}{\widehat\sigma_t}}_{w_{2t}}
  + \underbrace{\frac{\varepsilon_t}{\widehat\sigma_t}}_{\varepsilon_t^*} ,
  \qquad
  \D(\varepsilon_t^*) = \frac{\sigma_t^2}{\widehat\sigma_t^2} \approx 1 .
\]
Здесь $y_t^*$, $w_{1t}$, $w_{2t}$ — новые переменные, которые считаются по
данным, $\varepsilon_t^*$ — новая ошибка. Её дисперсия приблизительно
постоянна (ровно 1 была бы при точно известных $\sigma_t$). Оцениваем обычным
МНК регрессию $y_t^*$ на $w_{1t}$ и $w_{2t}$ \emph{без свободного члена}:
столбец единиц исходной модели после деления превратился в $w_{1t}$, и его
коэффициент — это $\alpha$. Полученные коэффициенты при $w_{1t}$ и $w_{2t}$ —
и есть искомые оценки $\widehat\alpha$ и $\widehat\beta$ исходной модели;
пересчитывать их обратно не нужно.

\medskip
\textbf{Что делает деление.} Наблюдения с большим разбросом получают меньший
вес. Например, пусть $\widehat\mu = 1$, $\widehat\delta = 0{,}5$:
\begin{center}
\begin{tabular}{cccc}
\toprule
$x_t$ & $\widehat\sigma_t^2 = 1 + 0{,}5x_t$ & $\widehat\sigma_t$ & вес в сумме квадратов $1/\widehat\sigma_t^2$ \\
\midrule
2 & 2 & 1,41 & 0,50 \\
8 & 5 & 2,24 & 0,20 \\
\bottomrule
\end{tabular}
\end{center}
Наблюдение с $x_t = 8$ более шумное, поэтому влияет на оценки в 2,5 раза
слабее, чем наблюдение с $x_t = 2$. Обычный МНК дал бы им одинаковый вес.

\medskip
\textbf{Итог: три регрессии.}
\begin{center}
\small
\begin{tabular}{@{}clll@{}}
\toprule
 & Зависимая переменная & Регрессоры & Результат \\
\midrule
1 & $y_t$ & константа, $x_t$ & остатки $e_t$ \\
2 & $e_t^2$ & константа, $x_t$ & $\widehat\mu$, $\widehat\delta$ $\Rightarrow$ $\widehat\sigma_t$ \\
3 & $y_t/\widehat\sigma_t$ & $1/\widehat\sigma_t$, $x_t/\widehat\sigma_t$ (без константы) & $\widehat\alpha$, $\widehat\beta$ \\
\bottomrule
\end{tabular}
\end{center}
Регрессии 1--2 — этап 1, регрессия 3 — этап 2. Оценки $\widehat\alpha$,
$\widehat\beta$ состоятельны и при большом $n$ эффективнее обычного МНК, а
стандартные ошибки регрессии 3 приближённо верны.
\end{solution}

% -------------------------------------------------------------------
\ifsolutions\clearpage\fi
\problem{6.6}

Рассмотрим модель, связывающую количество вакансий $w_t$ и уровень
безработицы $u_t$:
\[
  \ln w_t = \beta_1 + \beta_2\ln u_t + \varepsilon_t .
\]
Ошибки $\varepsilon_t$ независимы и нормально распределены
$N(0, \sigma_{\varepsilon}^2)$.
\begin{enumerate}[label=\asbuk*), leftmargin=2.2em, itemsep=3pt]
  \item Используя (искусственные) данные из таблицы, найдите МНК-оценки
        параметров $\beta_1$ и $\beta_2$, а также 95\,\%-ный доверительный
        интервал для $\beta_2$.
  \item Вычислите статистику Дарбина--Уотсона. Что её значение говорит об
        исходном предположении об ошибках $\varepsilon_t$? Что можно сказать
        о доверительном интервале, найденном в п.~а)?
  \item Оцените модель заново, используя модель автокорреляции первого
        порядка для ошибок регрессии. Найдите 95\,\%-ный доверительный
        интервал для $\beta_2$. Сравните результат с интервалом, полученным в
        п.~а).
\end{enumerate}

\begin{center}
\small
\begin{tabular}{rcc @{\qquad} rcc @{\qquad} rcc}
\toprule
$t$ & $w_t$ & $u_t$ & $t$ & $w_t$ & $u_t$ & $t$ & $w_t$ & $u_t$ \\
\midrule
1 & 1,73 & 8,65 &  9 & 5,06 & 2,87 & 17 & 3,15 & 4,72 \\
2 & 1,94 & 4,82 & 10 & 2,81 & 5,29 & 18 & 1,92 & 7,45 \\
3 & 3,05 & 2,67 & 11 & 4,43 & 3,31 & 19 & 2,26 & 6,21 \\
4 & 4,17 & 2,67 & 12 & 3,19 & 5,44 & 20 & 6,18 & 2,64 \\
5 & 2,52 & 2,58 & 13 & 2,23 & 6,80 & 21 & 2,07 & 8,55 \\
6 & 1,71 & 8,07 & 14 & 2,06 & 8,25 & 22 & 8,39 & 2,60 \\
7 & 1,95 & 8,83 & 15 & 3,33 & 3,44 & 23 & 2,75 & 6,25 \\
8 & 2,57 & 5,54 & 16 & 2,12 & 7,80 & 24 & 6,10 & 2,70 \\
\bottomrule
\end{tabular}
\end{center}

\begin{excel}
Добавьте столбцы \xl{=LN(w)} и \xl{=LN(u)} и оцените регрессию
\ui{Анализ данных} $\to$ \ui{Регрессия} (галочка \ui{Остатки}). DW по
остаткам $e_1, \dots, e_{24}$: \xl{=СУММКВРАЗН($e_2{:}e_{24}$; $e_1{:}e_{23}$)/СУММКВ($e_1{:}e_{24}$)}.
Для п.~в) — столбцы $y_t - \widehat{\rho}y_{t-1}$ и
$x_t - \widehat{\rho}x_{t-1}$ и регрессия первого на второй с константой. Это один шаг коррекции; для следующего шага пересчитайте остатки исходного уравнения.
\end{excel}

\begin{solution}
Обозначим $y_t=\ln w_t$, $x_t=\ln u_t$; $n=24$. Данные — временной ряд:
статистика DW сравнивает соседние по времени остатки, поэтому наблюдения
не сортируем. План: оценить модель и интервал (п.~а), проверить, независимы ли
ошибки (п.~б), и пересчитать оценки и интервал с учётом автокорреляции (п.~в).

\medskip
\textbf{а) МНК.}
\begin{center}
\begin{tabular}{lrrrr}
\toprule
Переменная & Коэффициент & Ст. ошибка & $t$-статистика & $P$-значение \\
\midrule
const   & $\phantom{-}2{,}2997$ & 0,1860 & $12{,}366$ & $<0{,}0001$ \\
$\ln u$ & $-0{,}7791$ & 0,1133 & $-6{,}875$ & $<0{,}0001$ \\
\midrule
\multicolumn{5}{l}{$R^2 = 0{,}6824$;\quad $\ESS = 1{,}4119$;\quad $\mathrm{DW} = 1{,}1026$} \\
\bottomrule
\end{tabular}
\end{center}
$\widehat{\beta}_1 = 2{,}30$, $\widehat{\beta}_2 = -0{,}78$: рост безработицы
на 1\,\% сопровождается в среднем снижением числа вакансий на 0,78\,\%.
Доверительный интервал по обычной формуле:
\[
  \widehat{\beta}_2 \pm t_{0{,}025}(n-2)\,s_{\widehat{\beta}_2}
  = -0{,}7791 \pm 2{,}074\cdot 0{,}1133 = (-1{,}014;\ -0{,}544),
\]
где $s_{\widehat\beta_2} = 0{,}1133$ — стандартная ошибка наклона из
таблицы, $t_{0{,}025}(22) = 2{,}074$ — критическое значение при
$n - 2 = 22$ степенях свободы.

\medskip
\textbf{б) Статистика Дарбина--Уотсона.} По остаткам $e_t$ \eqref{eq:dw}:
\[
  \mathrm{DW} = \frac{\sum_{t=2}^{24}(e_t - e_{t-1})^2}{\sum_{t=1}^{24}e_t^2}
              = \frac{1{,}5568}{1{,}4119} = 1{,}103,
  \qquad
  \widehat{\rho} \approx 1 - \frac{\mathrm{DW}}{2} = 0{,}45 ,
\]
где оценка $\widehat\rho$ получена из приближения
$\mathrm{DW} \approx 2(1 - \widehat\rho)$ в \eqref{eq:dw}. По таблице DW
для $n = 24$ и одного регрессора на уровне 5\,\% $d_L = 1{,}27$,
$d_U = 1{,}45$. $\mathrm{DW} = 1{,}10 < d_L$, поэтому гипотеза об отсутствии
автокорреляции отвергается: есть положительная автокорреляция первого
порядка. Значит, исходное предположение о \emph{независимости} ошибок
нарушено. Оценки МНК остаются несмещёнными, но стандартные ошибки из п.~а)
посчитаны по формулам, верным только для независимых ошибок, и поэтому
ненадёжны. Следовательно, интервал из п.~а) тоже ненадёжен: его настоящий
уровень доверия может заметно отличаться от 95\,\%.

\dwvisual

\medskip
\textbf{в) Оценка с учётом AR(1).} В п.~б) выяснилось, что соседние ошибки
связаны. Опишем эту связь моделью AR(1):
\[
  \varepsilon_t = \rho\varepsilon_{t-1} + \eta_t ,
\]
где $\rho$ — коэффициент автокорреляции ($|\rho| < 1$), $\eta_t$ —
независимые ошибки с постоянной дисперсией. Идея: преобразовать уравнение
так, чтобы в нём осталась только ошибка $\eta_t$. Тогда предпосылки МНК
выполнены, и обычные стандартные ошибки и интервал снова верны.

\smallskip
\emph{Шаг 1. Преобразование (пока считаем $\rho$ известным).} Запишем модель
для моментов $t$ и $t-1$, второе уравнение умножим на $\rho$:
\[
  y_t = \beta_1 + \beta_2x_t + \varepsilon_t ,
  \qquad
  \rho y_{t-1} = \rho\beta_1 + \rho\beta_2x_{t-1} + \rho\varepsilon_{t-1} .
\]
Вычтем второе из первого. Ошибка становится
$\varepsilon_t - \rho\varepsilon_{t-1} = \eta_t$, и получаем модель в
квазиразностях \eqref{eq:co}:
\[
  {\underbrace{y_t - \rho y_{t-1}}_{y_t^*}}
  = {\underbrace{\beta_1(1 - \rho)}_{\beta_1^*}}
  + \beta_2{\underbrace{(x_t - \rho x_{t-1})}_{x_t^*}} + \eta_t,
  \qquad t = 2, \dots, 24 ,
\]
где $y_t^*$, $x_t^*$ — новые переменные, $\beta_1^*$ — свободный член
нового уравнения. Это обычная регрессия $y_t^*$ на $x_t^*$ с константой.
Её наклон — это сразу $\beta_2$, а $\beta_1$ находится из константы:
$\beta_1 = \beta_1^*/(1 - \rho)$. Наблюдение $t = 1$ теряется: для него нет
$y_0$ и $x_0$, поэтому остаётся 23 наблюдения.

\smallskip
\emph{Шаг 2. Оценка $\rho$.} На деле $\rho$ неизвестно. Уравнение
$\varepsilon_t = \rho\varepsilon_{t-1} + \eta_t$ — регрессия $\varepsilon_t$
на $\varepsilon_{t-1}$ без константы. Заменяем ошибки остатками $e_t$ из
п.~а) и получаем МНК-оценку \eqref{eq:dw}:
\[
  \widehat\rho = \frac{\sum_{t=2}^{24} e_te_{t-1}}{\sum_{t=2}^{24} e_{t-1}^2}
  = \frac{0{,}5911}{1{,}3321} = 0{,}4437 .
\]
Это близко к грубой оценке $1 - \mathrm{DW}/2 = 0{,}45$ из п.~б).

\smallskip
\emph{Шаг 3. МНК для квазиразностей с $\widehat\rho = 0{,}4437$.} Для
каждого $t = 2, \dots, 24$ считаем новые переменные, например для $t = 2$:
\begin{align*}
  y_2^* &= y_2 - 0{,}4437\,y_1 = 0{,}6627 - 0{,}4437\cdot 0{,}5481 = 0{,}4195 ,\\
  x_2^* &= x_2 - 0{,}4437\,x_1 = 1{,}5728 - 0{,}4437\cdot 2{,}1576 = 0{,}6154 .
\end{align*}
Регрессия $y_t^*$ на $x_t^*$ по 23 наблюдениям даёт
$\widehat\beta_1^* = 1{,}3302$ и $\widehat\beta_2 = -0{,}8296$, откуда
$\widehat\beta_1 = 1{,}3302/(1 - 0{,}4437) = 2{,}3913$.

\smallskip
\emph{Шаг 4. Итерации Кохрейна--Оркатта.} Оценка $\widehat\rho$ из шага~2
построена по остаткам обычного МНК, который автокорреляцию не учитывал.
Теперь у нас есть уточнённые $\widehat\beta_1$, $\widehat\beta_2$. По ним
считаем новые остатки \emph{исходного} уравнения
$e_t = y_t - \widehat\beta_1 - \widehat\beta_2x_t$, заново оцениваем
$\widehat\rho$ (шаг~2) и повторяем шаг~3. Останавливаемся, когда
$\widehat\rho$ перестаёт меняться:
\begin{center}
\begin{tabular}{ccccc}
\toprule
Итерация & $\widehat\rho$ & $\widehat\beta_1^*$ & $\widehat\beta_1 = \widehat\beta_1^*/(1-\widehat\rho)$ & $\widehat\beta_2$ \\
\midrule
1 & 0,4437 & 1,3302 & 2,3913 & $-0{,}8296$ \\
2 & 0,4668 & 1,2768 & 2,3948 & $-0{,}8312$ \\
3 & 0,4675 & 1,2752 & 2,3949 & $-0{,}8312$ \\
4 & 0,4676 & 1,2751 & 2,3949 & $-0{,}8312$ \\
\bottomrule
\end{tabular}
\end{center}

\smallskip
\emph{Шаг 5. Итоговые оценки.} Последняя регрессия:
$y_t - 0{,}4676y_{t-1}$ на $x_t - 0{,}4676x_{t-1}$. Её константа
$\widehat\beta_1^* = 1{,}2751$ со стандартной ошибкой $0{,}0889$. Оба числа
делим на $1 - \widehat\rho = 0{,}5324$, чтобы вернуться к $\beta_1$:
$\widehat\beta_1 = 2{,}3949$, $s_{\widehat\beta_1} = 0{,}0889/0{,}5324 =
0{,}1669$. Наклон и его ошибка пересчёта не требуют. Стандартные ошибки
посчитаны так, будто $\widehat\rho$ известно точно:
\begin{center}
\begin{tabular}{lrrrr}
\toprule
Переменная & Коэффициент & Ст. ошибка & $t$-статистика & $P$-значение \\
\midrule
const   & $\phantom{-}2{,}3949$ & 0,1669 & $14{,}349$ & $<0{,}0001$ \\
$\ln u$ & $-0{,}8312$ & 0,0930 & $-8{,}939$ & $<0{,}0001$ \\
\midrule
\multicolumn{5}{l}{$\widehat\rho=0{,}4676$;\quad $\mathrm{DW} = 2{,}1617$} \\
\bottomrule
\end{tabular}
\end{center}

\noindent\begin{minipage}{\linewidth}
\emph{Шаг 6. Интервал.} Приближённый интервал для $\beta_2$ по конечной регрессии
(23 наблюдения, константа и наклон: $23-2=21$ степень свободы,
$t_{0{,}025}(21)=2{,}080$):
\[
  -0{,}8312 \pm 2{,}080\cdot 0{,}0930 = (-1{,}025;\ -0{,}638).
\]
Это приближение: дополнительная неопределённость оценки $\rho$
в обычной формуле стандартной ошибки последнего шага не учитывается.
\end{minipage}

\medskip
\noindent\begin{minipage}{\linewidth}
\textbf{Сравнение.}
\intervalvisual
Оценка наклона изменилась мало ($-0{,}83$ против $-0{,}78$): рост
безработицы на 1\,\% связан со снижением числа вакансий примерно на
0,83\,\%. Интервал с учётом AR(1) уже (длина 0,39 против 0,47). Но это не
доказывает, что точность выросла: интервал из п.~а) построен при нарушенной
предпосылке, и его ширине нельзя доверять. Второй интервал верен
приближённо и только если ошибки действительно следуют AR(1). Проверка:
у остатков последней регрессии $\widehat\eta_t$ статистика
$\mathrm{DW} = 2{,}16$, близко к 2, — автокорреляции первого порядка в них
не видно (графики ниже).
\end{minipage}
\arresidualvisual
\end{solution}

% -------------------------------------------------------------------
\ifsolutions\clearpage\fi
\problem{6.7}

В таблице представлены данные о потребительских расходах $C$ и располагаемом
доходе $Y_d$ тридцати семей (долл.): для каждого уровня дохода — расходы
трёх семей.
\begin{center}
\begin{tabular}{ccc|c}
\toprule
\multicolumn{3}{c|}{Потребление $C$} & Доход $Y_d$ \\
\midrule
10700 & 10900 & 11200 & 12000 \\
11400 & 11700 & 12100 & 13000 \\
12300 & 12600 & 13200 & 14000 \\
13000 & 13300 & 13600 & 15000 \\
13800 & 14000 & 14200 & 16000 \\
14400 & 14900 & 15300 & 17000 \\
15000 & 15700 & 16400 & 18000 \\
15900 & 16500 & 16900 & 19000 \\
16900 & 17500 & 18100 & 20000 \\
17200 & 17800 & 18500 & 21000 \\
\bottomrule
\end{tabular}
\end{center}
\begin{enumerate}[label=\asbuk*), leftmargin=2.2em, itemsep=3pt]
  \item Проведите регрессию $C$ на $Y_d$ и проверьте наличие или отсутствие
        гетероскедастичности.
  \item Если в п.~а) выявлена гетероскедастичность, осуществите коррекцию на
        гетероскедастичность.
\end{enumerate}

\begin{excel}
Данные — в «Семинар 6 (приложение).xlsx» (30 строк: $Y$, $C$). Регрессия —
\ui{Анализ данных} $\to$ \ui{Регрессия} с галочкой \ui{Остатки}. Для тестов
добавьте столбцы $e^2$, $|e|$, $Y^2$; Голдфелд--Квандт — две регрессии по
строкам 1--12 и 19--30 (данные уже упорядочены по $Y$); ранги —
\xl{=РАНГ.СР(\dots)}. Для взвешенного МНК — столбцы $C/Y$ и $1/Y$ и регрессия
$C/Y$ на $1/Y$ \emph{с константой}. Тридцать семей — тридцать
строк: каждый доход повторяется трижды; усреднять потребление по тройкам нельзя.
\end{excel}

\begin{solution}
Обозначим доход $Y$ ($= Y_d$); $n = 30$.

\medskip
\textbf{а) Регрессия} $C_t = \alpha + \beta Y_t + \varepsilon_t$:
\begin{center}
\begin{tabular}{lrrrr}
\toprule
Переменная & Коэффициент & Ст. ошибка & $t$-статистика & $P$-значение \\
\midrule
const & 1580,0 & 447,56 & 3,530 & 0,0015 \\
$Y$   & 0,7830 & 0,0267 & 29,30 & $<0{,}0001$ \\
\midrule
\multicolumn{5}{l}{$R^2 = 0{,}9684$;\quad $\ESS = 4\,948\,727$} \\
\bottomrule
\end{tabular}
\end{center}
Оценённая модель: $\widehat C_t=1580+0{,}78303Y_t$.
Дополнительный доллар дохода связан в среднем с 0,78 доллара потребления.
Остатки этой модели используем во всех тестах ниже. Например, для первой
семьи $\widehat C_1=10976{,}36$, $e_1=10700-10976{,}36=-276{,}36$,
$e_1^2\approx76377$. Стандартные ошибки в таблице верны, только если
дисперсия ошибок постоянна; высокий $R^2$ этого не гарантирует. Поэтому
сначала проверяем гетероскедастичность.

\medskip
\textbf{Проверка на гетероскедастичность.}

\smallskip
\noindent\begin{minipage}{\linewidth}
\emph{1) График.} Квадраты остатков против дохода:
\begin{center}
\begin{tikzpicture}[x=0.85cm, y=0.57cm]
  % x = Y/1000 - 11; y = e^2 / 100000
  \draw[->] (0,0) -- (10.8,0) node[anchor=north west] {$Y$};
  \draw[->] (0,0) -- (0,8.0) node[anchor=east] {$e_t^2$};
  \foreach \x/\l in {1/12000, 3/14000, 5/16000, 7/18000, 9/20000}
    \draw (\x,0) -- (\x,-0.12) node[anchor=north, font=\small] {\l};
  \foreach \y/\l in {2/200\,000, 4/400\,000, 6/600\,000}
    \draw (0,\y) -- (-0.12,\y) node[anchor=east, font=\small] {\l};
  \foreach \p in {(1,0.7638),(1,0.0583),(1,0.5001),(2,1.2916),(2,0.0353),(2,1.1601),(3,0.5877),(3,0.0331),(3,4.3241),(4,1.0592),(4,0.0065),(4,0.7538),(5,0.9516),(5,0.1177),(5,0.0838),(6,2.4159),(6,0.0007),(6,1.6686),(7,4.5501),(7,0.0065),(7,5.2628),(8,3.1089),(8,0.0180),(8,1.9574),(9,1.1601),(9,0.6729),(9,7.3856),(10,6.7838),(10,0.5001),(10,2.2692)}
    \fill[accent] \p circle (2.2pt);
\end{tikzpicture}

\smallskip
{\small\itshape С ростом дохода разброс остатков увеличивается: большие $e_t^2$
встречаются в основном при высоком доходе.\par}
\end{center}
\end{minipage}
График наводит на мысль о гетероскедастичности: дисперсия ошибок растёт
с доходом. Проверим это формально.

\smallskip
\emph{2) Тест Голдфелда--Квандта.} Наблюдения уже упорядочены по возрастанию
$Y$. Выбрасываем $c = 6$ средних семей (доход 16\,000 и 17\,000), остаются
$30 - 6 = 24$ семьи, по $m = 12$ в группе. Оцениваем регрессию $C$ на $Y$
отдельно по первым 12 и по последним 12 семьям и берём их остаточные суммы
квадратов: $\ESS_1 = 1\,046\,000$, $\ESS_2 = 3\,344\,000$. В каждой регрессии
$k = 2$ параметра (константа и наклон), поэтому $m - k = 10$ степеней
свободы. По \eqref{eq:gq}
\[
  F = \frac{\ESS_2/(m-k)}{\ESS_1/(m-k)} = \frac{\ESS_2/10}{\ESS_1/10},
\]
\[
  F = \frac{3\,344\,000/10}{1\,046\,000/10} = 3{,}20 > 2{,}98 = F_{0{,}05}(10, 10),
  \qquad P = 0{,}040 .
\]
Гипотеза гомоскедастичности отвергается: в группе с высоким доходом
оценка дисперсии ошибок примерно в 3,2 раза выше. Средние шесть семей
исключаем только для этого теста; остальные расчёты используют все 30
наблюдений.

\gqvisual

\smallskip
\emph{3) Тест Бройша--Пагана.} Идём по шагам из теории; дисперсия может
зависеть от дохода, поэтому $z_t = Y_t$ и $p = 1$.
\begin{enumerate}[label=\asbuk*), leftmargin=1.8em, itemsep=1pt]
  \item Остатки $e_t$ исходной модели уже есть (п.~а).
  \item $\widehat{\sigma}^2 = \ESS/n = 4\,948\,727/30 = 164\,958$ (делим на
        $n = 30$, а не на $n - k = 28$).
  \item $g_t = e_t^2/164\,958$; например, для первой семьи
        $g_1 = 76\,377/164\,958 = 0{,}463$.
  \item Регрессия $g_t = a_0 + a_1Y_t + v_t$; её объяснённая сумма квадратов
        (в Excel — строка «Регрессия» таблицы дисперсионного анализа)
        $\RSS_{\text{aux}} = 8{,}22$.
  \item По \eqref{eq:bp}
        \[
          \mathrm{BP} = \frac{\RSS_{\text{aux}}}{2} = \frac{8{,}22}{2} = 4{,}11
          > 3{,}84 = \chi^2_{0{,}05}(1),
          \qquad P = 0{,}043 .
        \]
\end{enumerate}
Гомоскедастичность снова отвергается. Вариант Кёнкера, не требующий
нормальности ошибок: $R^2$ той же вспомогательной регрессии равен $0{,}183$,
и $\mathrm{BP}_K = nR^2_{\text{aux}} = 30\cdot 0{,}183 = 5{,}49 > 3{,}84$,
$P = 0{,}019$ — вывод тот же.

\smallskip
\emph{4) Тест Уайта.} Регрессор в модели один, поэтому произведений нет, и
вспомогательная регрессия содержит $Y_t$ и $Y_t^2$:
\[
  e_t^2 = a_0 + a_1Y_t + a_2Y_t^2 + v_t ,
  \qquad
  H_0\colon a_1 = a_2 = 0 .
\]
Её коэффициент детерминации $R^2_{\text{aux}} = 0{,}194$; проверяемых
коэффициентов два, поэтому $p = 2$. По \eqref{eq:white}
\[
  W = nR^2_{\text{aux}} = 30\cdot 0{,}194 = 5{,}82 < 5{,}99 = \chi^2_{0{,}05}(2),
  \qquad P = 0{,}055 .
\]
На выбранном заранее уровне 5\,\% $H_0$ не отвергаем (на 10\,\% отвергли
бы). Расхождение с тестом Бройша--Пагана объяснимо: у Уайта проверяются два
коэффициента вместо одного, а лишняя степень свободы при $n = 30$ снижает
мощность теста.

\smallskip
\emph{5) Тест Глейзера.} Модуль остатка $|e_t|$ показывает, насколько
семья отклонилась от линии регрессии, и поэтому служит грубой оценкой
стандартного отклонения $\sigma_t$ её ошибки. Например, у первой семьи
$|e_1| = 276{,}36$. Если $\sigma_t$ растёт с доходом, то и $|e_t|$ в среднем
растут с $Y_t$. Это проверяют регрессией
\[
  |e_t| = a_0 + a_1Y_t^{\gamma} + v_t ,
\]
где $a_0$, $a_1$ — параметры, $v_t$ — ошибка, $\gamma$ — показатель, который
задаёт форму зависимости. Какая форма верна, заранее неизвестно, поэтому
пробуем три: $\gamma = 1/2$, $1$, $2$ ($\sigma_t$ растёт как $\sqrt{Y_t}$,
$Y_t$ или $Y_t^2$). Для каждой формы проверяем
\[
  H_0\colon a_1 = 0 \ (\text{$|e_t|$ не зависит от дохода}),
  \qquad
  t = \frac{\widehat a_1}{s_{\widehat a_1}} ,
\]
где $s_{\widehat a_1}$ — стандартная ошибка наклона. $H_0$ отвергаем, если
$|t| > t_{0{,}025}(28) = 2{,}048$; $28 = n - 2$ степеней свободы.

\smallskip
\emph{Расчёт для $\gamma = 1$.} Регрессия $|e_t|$ на константу и $Y_t$ даёт
$\widehat a_0 = -228{,}8$, $\widehat a_1 = 0{,}03364$,
$s_{\widehat a_1} = 0{,}01459$:
\[
  t = \frac{0{,}03364}{0{,}01459} = 2{,}305 > 2{,}048,
  \qquad P = 0{,}029 .
\]
Значит, $H_0$ отвергается. Наклон означает: при росте дохода на 1000 долл.
средний модуль остатка растёт примерно на 34 долл., от 175 при
$Y = 12\,000$ до 478 при $Y = 21\,000$.

\smallskip
\noindent\begin{minipage}{\linewidth}
\parindent=1.5em
\emph{Все три формы:}
\begin{center}
\begin{tabular}{lcccc}
\toprule
Регрессия & $\widehat a_1$ & $t$ наклона & $P$ (прибл.) & $R^2$ \\
\midrule
$|e_t|$ на $\sqrt{Y_t}$ & 8,467 & 2,270 & 0,031 & 0,155 \\
$|e_t|$ на $Y_t$ & 0,03364 & 2,305 & 0,029 & 0,159 \\
$|e_t|$ на $Y_t^2$ & $1{,}037\cdot 10^{-6}$ & 2,361 & 0,025 & 0,166 \\
\bottomrule
\end{tabular}
\end{center}
\end{minipage}
Во всех трёх регрессиях $t > 2{,}048$, наклон положителен: разброс ошибок
растёт с доходом. $P$-значения приближённые, потому что вместо ошибок
$\varepsilon_t$ использованы остатки $e_t$.

Выбрать одну форму по этим данным нельзя: $R^2$ почти одинаковы
(0,155--0,166). На отрезке дохода от 12\,000 до 21\,000 величины $\sqrt{Y}$,
$Y$ и $Y^2$ почти линейно связаны друг с другом (попарные корреляции
больше 0,99), поэтому три регрессии дают по сути один и тот же результат,
а не три независимых доказательства. Для коррекции в п.~б) берём самую
простую форму $\sigma_t \propto Y_t$.

\smallskip
\emph{6) Ранговая корреляция Спирмена} между $|e_t|$ и $Y_t$. Если разброс
ошибок растёт с доходом, то семьи с бо́льшим доходом в среднем имеют бо́льшие
$|e_t|$. Тест Спирмена проверяет это по рангам: форма зависимости не важна,
важно только, растёт ли $|e_t|$ вместе с $Y_t$.

\smallskip
\emph{Шаг 1. Ранги дохода $R_Y$.} Ранг — номер значения в ряду,
упорядоченном по возрастанию. Каждый уровень дохода встречается трижды (в данных для каждого дохода
приведены расходы трёх семей), и
одинаковым значениям даём средний ранг занимаемых ими мест. Доход 12\,000
занимает места 1, 2, 3, поэтому у всех трёх семей ранг $(1+2+3)/3 = 2$;
доход 13\,000 занимает места 4--6 и получает ранг 5; и так далее до
21\,000 с рангом 29.

\smallskip
\emph{Шаг 2. Ранги модулей остатков $R_{|e|}$.} Упорядочиваем 30 значений
$|e_t|$ по возрастанию и нумеруем места:

\spearmanrankvisual
Наименьший модуль $8{,}48$ (семья 17) получает ранг 1. Модуль $25{,}45$
встречается дважды (семьи 11 и 20); они занимают места 2 и 3 и получают
ранг $(2+3)/2 = 2{,}5$. Так же совпадают $223{,}64$ (семьи 3 и 29, места
10--11, ранг 10,5) и $340{,}61$ (семьи 6 и 25, места 18--19, ранг 18,5).

\smallskip
\noindent\begin{minipage}{\linewidth}
\parindent=1.5em
\emph{Шаг 3. Разности рангов.} Возвращаем семьи в порядок роста дохода
и для каждой считаем $d_t = R_Y - R_{|e|}$:

\spearmanvisual
\end{minipage}
Например, у семьи 1: $R_Y = 2$, $R_{|e|} = 15$, $d_1 = -13$, $d_1^2 = 169$.

\smallskip
\emph{Шаг 4. Коэффициент Спирмена.} По \eqref{eq:spearman}
\[
  r_S = 1 - \frac{6\sum d_t^2}{n(n^2-1)} = 1 - \frac{6\cdot 2931{,}5}{30\cdot 899}
      = 1 - 0{,}652 = 0{,}348 .
\]
Из-за совпадающих рангов эта формула приближённая. Точное значение —
обычный коэффициент корреляции столбцов $R_Y$ и $R_{|e|}$ (в Excel
\xl{=КОРРЕЛ(\dots)}): $r_S = 0{,}345$. Разница мала, дальше используем
$0{,}345$. Знак положительный: при большем доходе модули остатков в среднем
больше.

\smallskip
\emph{Шаг 5. Значимость.} $H_0$: ранговой связи нет. По \eqref{eq:spearman}
\[
  t = \frac{r_S\sqrt{n-2}}{\sqrt{1 - r_S^2}}
    = \frac{0{,}345\sqrt{28}}{\sqrt{1 - 0{,}345^2}} = 1{,}94 < 2{,}048 = t_{0{,}025}(28),
  \qquad P \approx 0{,}062 .
\]
На уровне 5\,\% $H_0$ не отвергаем, на уровне 10\,\% отвергли бы.

\smallskip
\textbf{Сводка диагностики при уровне 5\,\%.}
\begin{center}
\begin{tabular}{lcl}
\toprule
Метод & $P$ & Решение о нулевой гипотезе \\
\midrule
Голдфелд--Квандт & 0,040 & Отвергаем \\
Бройш--Паган & 0,043 & Отвергаем \\
Уайт & 0,055 & Не отвергаем \\
Глейзер: $|e|$ на $Y$ & 0,029 & Отвергаем нулевой наклон \\
Спирмен & $\approx0{,}062$ & Не отвергаем нулевую корреляцию \\
\bottomrule
\end{tabular}
\end{center}
График, Голдфелд--Квандт, Бройш--Паган и Глейзер указывают на рост
дисперсии с доходом. Уайт и Спирмен показывают ту же направленность, но
значимы только на уровне 10\,\%. Неотвержение $H_0$ не доказывает
гомоскедастичность: при $n = 30$ у этих тестов малая мощность. Вывод:
гетероскедастичность есть, и её нужно учесть.

\smallskip
\textbf{б) Коррекция: взвешенный МНК.} Для взвешенного МНК нужна форма
зависимости дисперсии от дохода. Её подсказывает тест Глейзера: $|e_t|$
— грубая мера стандартного отклонения ошибки, и регрессия $|e_t|$ на $Y_t$
даёт значимый положительный наклон, то есть масштаб ошибок растёт примерно
пропорционально доходу. Поэтому в качестве рабочей модели примем
$\sigma_t = \sigma Y_t$, где $\sigma$ — неизвестная постоянная (так же
сделано в сборнике), и по \eqref{eq:wls} с $z_t = Y_t$ разделим уравнение
на $Y_t$:
\[
  \frac{C_t}{Y_t} = \alpha\,\frac{1}{Y_t} + \beta + \frac{\varepsilon_t}{Y_t},
  \qquad
  \D\Bigl(\frac{\varepsilon_t}{Y_t}\Bigr) = \frac{\sigma^2Y_t^2}{Y_t^2} = \sigma^2 .
\]
Ошибка нового уравнения гомоскедастична. В нём $\beta$ стала свободным
членом, а $\alpha$ — коэффициентом при $1/Y_t$. Поэтому оцениваем регрессию
$C_t/Y_t$ на $1/Y_t$ с константой:
\begin{center}
\begin{tabular}{lrrrr}
\toprule
Переменная & Коэффициент & Ст. ошибка & $t$-статистика & $P$-значение \\
\midrule
$1/Y$ ($\alpha$) & 1557,0 & 392,28 & 3,969 & 0,0005 \\
const ($\beta$)  & 0,7844 & 0,0249 & 31,46 & $<0{,}0001$ \\
\midrule
\multicolumn{5}{l}{$R^2 = 0{,}360$ (для $C/Y$);\quad $R^2 = 0{,}968$ для исходной переменной $C$} \\
\bottomrule
\end{tabular}
\end{center}
Подставляя оценки в исходное уравнение, получаем
\[
  \widehat{C} = 1557{,}0 + 0{,}7844\,Y .
\]
\emph{Помогло ли деление?} Возьмём остатки преобразованной модели
$\widehat v_t=C_t/Y_t-\widehat\alpha/Y_t-\widehat\beta$ (а не остатки в
долларах) и проведём для них тест Уайта: регрессия $\widehat v_t^{\,2}$ на
константу, $1/Y_t$ и $1/Y_t^2$ даёт
\[
  nR^2_{\text{aux}} = 1{,}06 < 5{,}99 = \chi^2_{0{,}05}(2),
  \qquad P = 0{,}59 .
\]
Гомоскедастичность не отвергается: после деления на доход признаков
непостоянной дисперсии не видно. Это согласуется с выбранной моделью
$\sigma_t = \sigma Y_t$, хотя и не доказывает её.

\wlsresidualvisual

\begin{remark}
Наклоны МНК и ВМНК близки ($0{,}7830$ и $0{,}7844$): обе оценки
несмещённые, но в другой выборке они могли бы разойтись сильнее.
Стандартные ошибки ($0{,}0267$ у МНК и $0{,}0249$ у ВМНК) напрямую сравнивать
нельзя: первая неверна при гетероскедастичности, вторая верна, только если
верна модель $\sigma_t = \sigma Y_t$. $R^2$ преобразованной модели ($0{,}36$)
нельзя сравнивать с $R^2$ исходной ($0{,}968$): у них разные зависимые
переменные ($C/Y$ и $C$), поэтому снижение $R^2$ не означает, что модель
стала хуже.
\end{remark}
\end{solution}
